%% kapa: exam=1 round=33 year=2022 period=2 form=A subject=appraisal_law subject_ko=감정평가관계법규 seq=35 number=35
%% source: Q-Net 감정평가사 기출문제 (cst00309) artlSeq=5209997 file=raw/1008/5209997/2022년도 제33회 감정평가사 1차시험 2교시(감정평가 관계법규, 회계학).pdf
%% conversion: kordoc 4.24.51 -> tools/gen1.py (auto); review=proofread-sonnet; flags=none
\begin{question}{35}
부동산등기법령상 A(용익권 또는 담보권)와 B(등기원인에 그 약정이 있는 경우에만 기록하여야 하는 사항)의 연결로 옳지 않은 것은?
\begin{choices}
  \item A: 지역권, B: 범위
  \item A: 전세권, B: 존속기간
  \item A: 저당권, B: 변제기
  \item A: 근저당권, B: 존속기간
  \item A: 지상권, B: 지료와 지급시기
\end{choices}
\end{question}
\answer{1}
